\section{Related Work}
\label{sec:related}

\paragraph{Multi-view quality for instruction data.}
Instruction-data quality is inherently multi-faceted. Huang
et al.~\cite{huang2024ivm} formalize instruction mining through three
views---diversity, complexity, and accuracy---and fuse LoRA-space
coverage, uncertainty-based complexity, and reward-model accuracy to
select a compact tuning subset. Closely related instruction-selection
methods likewise combine several notions of quality: DEITA
integrates complexity, quality, and diversity~\cite{deita}; MoDS
balances quality, coverage, and necessity~\cite{mods}; SelectIT uses
uncertainty-aware self-reflection~\cite{selectit}; and IFD/Q2Q ranks
examples by model-derived instruction-following
difficulty~\cite{q2q}. These works establish that no single scalar
fully describes useful instruction data. Their views, however, are
typically evaluated and fused before fine-tuning. Our setting adds a
temporal distinction: pair reliability is assessed once, whereas the
usefulness of a reliable pair is refreshed as the model learns.

\paragraph{Quality filtering under unreliable supervision.}
Real instruction pools contain noisy, mismatched, incomplete, or
otherwise unreliable supervision. Quality-filtering methods address
this problem using learned estimators or stronger LLM
judges~\cite{alpagasus,cpqs2026}, while reward-model accuracy in
multi-view instruction mining provides an explicit correctness
signal~\cite{huang2024ivm}. More broadly, trusted multi-view learning
studies how evidence sources with unequal or conflicting reliability
should influence fusion~\cite{tmc2021,rcml2024}. TAG brings this
reliability-aware perspective to instruction-data selection. Its
static view is an instruction--response dependency detector whose
continuous gate attenuates questionable supervision before two
model-dependent utility signals are combined. This role differs from
both a global quality ranking and a corruption-type classifier: the
gate estimates the credibility of the pair, while dynamic selection
decides how useful that pair is at the current training stage.

\paragraph{Static target-conditioned and representation-based
selection.}
Other selectors define usefulness relative to a target distribution
or internal model geometry. DSIR uses importance resampling for
language-model data selection~\cite{dsir}; LESS scores candidates by
gradient similarity to a validation or few-shot target
set~\cite{less}; and Approx-LESS improves the scalability of this
influence calculation~\cite{approxless}. At the representation level,
NAIT extracts target-derived directions from in-domain seed examples
and ranks candidates by projecting their activations onto those
directions~\cite{nait2026}. This idea is related to evidence that
linear activation directions encode meaningful concepts or
capabilities, including TCAV~\cite{tcav}, representation
engineering~\cite{repe}, and contrastive activation
addition~\cite{caa}. These methods expose useful target or geometric
structure, but the target examples and representation directions are
generally fixed before fine-tuning. They therefore do not adapt their
geometric readout to later model states.

\paragraph{Training-adaptive data selection.}
Dynamic selectors address the complementary fact that example utility
changes during training. Active instruction tuning uses model
feedback to prioritize prompt-sensitive tasks~\cite{kung2023activeit},
and difficulty--uncertainty selectors refresh scalar rewards under
the current checkpoint. Data Agent formulates the process as a
sequential decision problem and learns a PPO-based selection
policy~\cite{dataagent}. In reinforcement fine-tuning, DOTS performs
online question selection using rollout-dependent reward
statistics~\cite{sun2025dots}; its verifier rewards and rollout groups
differ from the supervised instruction--response setting considered
here. These approaches make selection responsive to the evolving
model, but their explicit selection signal remains a scalar reward or
a learned policy output rather than an interpretable online
representation direction.

\paragraph{Positioning of TAG.}
TAG combines the two distinctions emphasized above: multiple quality
views play different roles, and only model-dependent usefulness needs
to be refreshed throughout training. A static counterfactual view
assigns a continuous reliability gate to each
instruction--response pair. Among candidates admitted by that gate, a
difficulty--uncertainty reward measures current informativeness, and
a state-conditioned representation-anchor view assigns a normalized
projection onto principal directions extracted from the current
model's probe displacements. The resulting anchor modulation is
bounded and independently ablatable rather than a separate policy.
This organization connects multi-view instruction
mining~\cite{huang2024ivm} with training-adaptive selection while
explicitly addressing low-quality supervision. The accompanying
analysis is correspondingly role-specific: total-variance
decomposition motivates separating scalar reward from anchor
projection, and standard spectral perturbation tools characterize the
local stability of the recomputed principal direction without claiming
stability of candidate projections or of the full
ranking~\cite{davis1970rotation,stewart1990matrix}.



